\section{Evaluation Metrics}
\label{app:metrics}

The primary safety metrics are unsafe-action count, availability violations, data-loss events, unnecessary drain rate, and pre-failure protection rate. Effectiveness metrics include time at risk, risk-exposure area from \cref{eq:risk-exposure}, repair completion time, and bytes rewritten. Performance metrics include foreground throughput and P50/P95/P99/P99.9 latency together with maintenance progress.

\section{Planned Fault and Workload Matrix}
\label{app:matrix}

Fault traces cover transient risk spikes, sustained risk, false-positive recovery, hard EIO, silent corruption, multiple concurrent risky devices, insufficient redundancy, metadata conflicts, and device addition followed by rebalance. Foreground workloads cover sequential reads/writes, random reads/writes, mixed read/write traffic, and metadata-heavy small-file operations.

\section{Reproducibility Protocol}
\label{app:reproducibility}

\begin{itemize}[leftmargin=*]
  \item Every experiment records the \argosfs{} commit and paper-experiment configuration.
  \item Randomized traces use explicit seeds that are replayed across all baselines.
  \item Raw JSONL/CSV records are retained before aggregation.
  \item Performance experiments report repetitions and confidence intervals.
  \item Device-mapper and QEMU artifacts retain fault parameters and execution logs.
  \item Published figures are generated from retained machine-readable results.
\end{itemize}
